% GENERATED FILE. Edit canonical lessons, then run npm run generate:books.
% canonical-source-hash: fnv1a64:3634f1d291c9ea37
% canonical-lessons: TE-C289-bhogi, TE-C289-panduga, TE-C289-jatara, TE-C289-veduka, TE-C289-suryodayam

\chapter{Bhogi, Festival, Village fair, Celebration, Sunrise}
\label{ch:te-bhogi-festival-village-fair-celebration-sunrise}
\hlchaptermodality{289}

\begin{chapteropening}
\textbf{By the end of this chapter:} \emph{I can say Bhogi, a festival, a village fair, a celebration and sunrise.}

Complete the last lesson of chapter 289: I can say Bhogi, a festival, a village fair, a celebration and sunrise.
\end{chapteropening}

\section[\texorpdfstring{\te{భోగి} (bhōgi)}{భోగి (bhōgi)}]{\te{భోగి} (bhōgi) — Bhogi, the first day of Sankranti}

\label{lesson:TE-C289-bhogi}



\begin{quote}
Before the new one: say the Telugu for diabetes, then the Telugu for blood pressure.
\end{quote}

\subsection*{You'll want to know: \te{భోగి}}
\textbf{\te{భోగి}} — \emph{bhōgi} — \textquotedblleft{}Bhogi, the first day of Sankranti\textquotedblright{}.

\begin{grammarlens}[title={What you've built}]
One more word for this chapter.
\end{grammarlens}

\subsection*{Your turn}
\begin{itemize}
\raggedright
  \item \emph{Say it:} \emph{bhōgi}
  \item \emph{Say it:} \emph{bhōgi}, once more
  \item \emph{Say it:} say \emph{bhōgi}
  \item \emph{Recall:} say the Telugu for diabetes, then the Telugu for blood pressure, then say \emph{bhōgi} again
\end{itemize}

\begin{tcolorbox}[breakable,colback=teal!4,colframe=teal!35!black,title={Before you move on}]
What does \te{భోగి} mean? (\textquotedblleft{}Bhogi, the first day of Sankranti\textquotedblright{}.) Say it once more.
\end{tcolorbox}
\section[\texorpdfstring{\te{పండుగ} (paṇḍuga)}{పండుగ (paṇḍuga)}]{\te{పండుగ} (paṇḍuga) — a festival}

\label{lesson:TE-C289-panduga}



\begin{quote}
Before the new one: say the Telugu for blood pressure, then the Telugu for Bhogi.
\end{quote}

\subsection*{You'll want to know: \te{పండుగ}}
\textbf{\te{పండుగ}} — \emph{paṇḍuga} — \textquotedblleft{}a festival\textquotedblright{}.

\begin{grammarlens}[title={What you've built}]
One more word for this chapter.
\end{grammarlens}

\subsection*{Your turn}
\begin{itemize}
\raggedright
  \item \emph{Say it:} \emph{paṇḍuga}
  \item \emph{Say it:} \emph{paṇḍuga}, once more
  \item \emph{Say it:} say \emph{paṇḍuga}
  \item \emph{Recall:} say the Telugu for blood pressure, then the Telugu for Bhogi, then say \emph{paṇḍuga} again
\end{itemize}

\begin{tcolorbox}[breakable,colback=teal!4,colframe=teal!35!black,title={Before you move on}]
What does \te{పండుగ} mean? (\textquotedblleft{}A festival\textquotedblright{}.) Say it once more.
\end{tcolorbox}
\section[\texorpdfstring{\te{జాతర} (jātara)}{జాతర (jātara)}]{\te{జాతర} (jātara) — a village fair, a temple festival}

\label{lesson:TE-C289-jatara}



\begin{quote}
Before the new one: say the Telugu for Bhogi, then the Telugu for a festival.
\end{quote}

\subsection*{You'll want to know: \te{జాతర}}
\textbf{\te{జాతర}} — \emph{jātara} — \textquotedblleft{}a village fair, a temple festival\textquotedblright{}.

\begin{grammarlens}[title={What you've built}]
One more word for this chapter.
\end{grammarlens}

\subsection*{Your turn}
\begin{itemize}
\raggedright
  \item \emph{Say it:} \emph{jātara}
  \item \emph{Say it:} \emph{jātara}, once more
  \item \emph{Say it:} say \emph{jātara}
  \item \emph{Recall:} say the Telugu for Bhogi, then the Telugu for a festival, then say \emph{jātara} again
\end{itemize}

\begin{tcolorbox}[breakable,colback=teal!4,colframe=teal!35!black,title={Before you move on}]
What does \te{జాతర} mean? (\textquotedblleft{}A village fair, a temple festival\textquotedblright{}.) Say it once more.
\end{tcolorbox}
\section[\texorpdfstring{\te{వేడుక} (vēḍuka)}{వేడుక (vēḍuka)}]{\te{వేడుక} (vēḍuka) — a celebration, a ceremony}

\label{lesson:TE-C289-veduka}



\begin{quote}
Before the new one: say the Telugu for a festival, then the Telugu for a village fair.
\end{quote}

\subsection*{You'll want to know: \te{వేడుక}}
\textbf{\te{వేడుక}} — \emph{vēḍuka} — \textquotedblleft{}a celebration, a ceremony\textquotedblright{}.

\begin{grammarlens}[title={What you've built}]
One more word for this chapter.
\end{grammarlens}

\subsection*{Your turn}
\begin{itemize}
\raggedright
  \item \emph{Say it:} \emph{vēḍuka}
  \item \emph{Say it:} \emph{vēḍuka}, once more
  \item \emph{Say it:} say \emph{vēḍuka}
  \item \emph{Recall:} say the Telugu for a festival, then the Telugu for a village fair, then say \emph{vēḍuka} again
\end{itemize}

\begin{tcolorbox}[breakable,colback=teal!4,colframe=teal!35!black,title={Before you move on}]
What does \te{వేడుక} mean? (\textquotedblleft{}A celebration, a ceremony\textquotedblright{}.) Say it once more.
\end{tcolorbox}
\section[\texorpdfstring{\te{సూర్యోదయం} (sūryōdayaṁ)}{సూర్యోదయం (sūryōdayaṁ)}]{\te{సూర్యోదయం} (sūryōdayaṁ) — sunrise}

\label{lesson:TE-C289-suryodayam}



\begin{quote}
Before the new one: say the Telugu for a village fair, then the Telugu for a celebration.
\end{quote}

\subsection*{You'll want to know: \te{సూర్యోదయం}}
\textbf{\te{సూర్యోదయం}} — \emph{sūryōdayaṁ} — \textquotedblleft{}sunrise\textquotedblright{}.

\begin{grammarlens}[title={What you've built}]
One more word for this chapter.
\end{grammarlens}

\subsection*{Your turn}
\begin{itemize}
\raggedright
  \item \emph{Say it:} \emph{sūryōdayaṁ}
  \item \emph{Say it:} \emph{sūryōdayaṁ}, once more
  \item \emph{Say it:} say \emph{sūryōdayaṁ}
  \item \emph{Recall:} say the Telugu for a village fair, then the Telugu for a celebration, then say \emph{sūryōdayaṁ} again
\end{itemize}

\begin{tcolorbox}[breakable,colback=teal!4,colframe=teal!35!black,title={Before you move on}]
What does \te{సూర్యోదయం} mean? (\textquotedblleft{}Sunrise\textquotedblright{}.) Say it once more.
\end{tcolorbox}
